\section{A variant with every derivative globally bounded}
\label{app:all-derivatives}

The activation in Theorem~\ref{thm:main} has global bounds on any
prescribed finite number of derivatives. Its sufficiently high
derivatives need not be bounded. This feature is not necessary for an
exponential dependence on teacher depth: a smaller encoding amplitude gives
one activation whose derivatives of \emph{every} order are globally
bounded. Amplifier weights of order $d$ preserve the linear teacher depth.

\begin{theorem}\label{thm:all-derivatives}
Fix a finite integer $r\ge1$ and $\eta>0$. There is one computable
activation $\phi$ with all the shape and $C^r$-closeness properties of
Theorem~\ref{thm:main}, such that every derivative of $\phi$ is globally
bounded. The same circuit-hard targets $f_d$ have globally unit-margin
teachers with
\[
 L_d=d+4,\qquad
 \text{active width }O(\log(d+2)),\qquad
 \text{coefficient magnitude }O(d+1),
\]
and polynomial coefficient bit length. For some $c>0$, every sample
of at most $\exp(cd)$ observations has posterior confidence-failure
rate at least $139/1024$, at every fixed Gaussian exponent
$\gamma\ge1$ and every sufficiently large width
$N\ge\exp(C_\gamma d)$. The fixed-width cold-limit conclusion holds
at every width above the memorizer threshold~\eqref{eq:MB}.

Thus this variant retains the sample lower bound
$\exp(\Omega(L_d))$. It does not assert a uniform
analytic-strip bound, a finite-order Gevrey bound, or efficient
activation evaluation.
\end{theorem}

\subsection{The modified bands and derivative estimates}

Fix a sufficiently large integer $M_0\ge24$, and set
\begin{align*}
 m_d&=M_0+\lceil\log_2(d+2)\rceil,&
 J_d&=2^{m_d+10},&
 \Delta_d&=2^{-d-4},\\
 r_d&=\Delta_d/J_d,&
 c_{d,k}&=-4d-\tfrac14-\Delta_d k,&
 \delta_d&=\tau r_d2^{-(m_d-3)d}.
\end{align*}
In particular,
\begin{equation}\label{eq:all-derivative-amplitude}
 \delta_d=\tau2^{-(m_d-2)d-m_d-14}.
\end{equation}
Use these parameters in~\eqref{eq:S}--\eqref{eq:phi}, with a fixed
sufficiently small positive rational $\tau$. The truth tables and
centers are otherwise unchanged.

For each fixed derivative order $j$, the Gaussian lattice estimate
in Lemma~\ref{lem:phi} gives
\begin{equation}\label{eq:all-derivative-bound}
 \sup_{\R}|S_d^{(j)}|
 \le C_j\tau(16J_d)^{j-1}2^{-(m_d-j-2)d}.
\end{equation}
Here $S_d$ denotes the whole $d$th band. This is summable in $d$
for every fixed $j$, because $m_d=\log_2d+O(1)$ and $J_d=O(d)$.
Consequently every derivative is globally bounded. On each complex
compact, the moving band centers and Gaussian suppression still
give locally uniform convergence; hence $S$ is entire. The explicit
tails and the computable finite truth tables give computability.

The weighted first- and second-derivative estimates have band
majorants
\[
 C\tau e^{8d}2^{-(m_d-3)d},\qquad
 C\tau J_de^{8d}2^{-(m_d-4)d},
\]
respectively. Both are summable. The cubic factor at zero supplies
the same extra powers of $u$ as in Lemma~\ref{lem:phi}. Therefore
\eqref{eq:weightedfirst}--\eqref{eq:weightedsecond} hold, with fixed
constants for this modified activation. A sufficiently small $\tau$
preserves oddness, slope one at zero, strict increase, global
Lipschitz constant one, positive-half-line strict concavity, and
saturation to $\pm1$. Equation~\eqref{eq:all-derivative-bound} makes
the prescribed $C^r$ distance smaller than $\eta$ after a further
fixed reduction of $\tau$. The constants bounding derivative $j$
can grow rapidly with $j$; no uniform analytic-type estimate is used.

\subsection{A growing finite-difference decoder}

Use $h_d=\Delta_d/(2m_d)$ and the $m_d$th finite difference in
\eqref{eq:decoder}. At $c_{d,k}$ the central bump contributes
$\delta_df_d(k)$. Every other query is at least $h_d$ from every
center in this band. Since $J_d=2^{m_d+10}$, its leakage, multiplied
by the coefficient sum $2^{m_d}=O(d)$, is $o(\delta_d)$.

For another band $e\ne d$, the logarithm of its amplitude ratio
is $O(d\log d+e\log e)$, whereas Gaussian suppression has exponent
at most $-c(d-e)^24^eJ_e^2$. Split the sum into $e\le d/2$,
$d/2<e<d$, and $e>d$. In the first range even the smallest fixed
$e$ gives a negative quadratic in $d$, dominating $d\log d$.
The second has suppression at least $\exp(-c4^{d/2})$; in the
third the amplitudes decrease and the Gaussian suppression is
stronger still. Polynomial factors and the coefficient sum do not
change these conclusions. Reflected bands are farther away.

For $m=m_d=O(\log d)$, the negative-half-line exponential series of
tanh gives
\[
 \sup_{u\le-4d+1}|\tanh^{(m)}u|
 \le C2^m e^{-8d+2}.
\]
Indeed the differentiated series is bounded by
$2^{m+1}\sum_{n\ge1}n^me^{-(8d-2)n}$, whose first term dominates
uniformly for large $d$ since $m=o(d)$. The tanh contribution to
the finite difference is at most
$C\Delta_d^m m^{-m}e^{-8d}$. Dividing by
\eqref{eq:all-derivative-amplitude} bounds its relative size by
\[
 C\tau^{-1}2^{-2d-3m+14}m^{-m}e^{-8d}=o(1).
\]
We therefore retain the uniform decoding estimate
$F_d(c_{d,k})=\delta_d(f_d(k)+o(1))$.

The first layer has $m_d+1$ neurons and evaluates these query
points. The next layer uses binomial coefficients of magnitude at
most $2^{m_d}=O(d)$. Set the amplifier gain
\[
 a_d=2^{m_d-2}=\Theta(d).
\]
Exactly $d+1$ scalar amplification layers give ideal final amplitude
\[
 \delta_da_d^{d+1}=\tau2^{-16}=: \rho.
\]
Their ideal preactivation amplitudes form a geometric sequence with
sum at most $2\rho$, since $a_d\ge2$. The estimate
\eqref{eq:taylor} therefore gives a uniform $O(\rho)$ relative-error
bootstrap, independent of the growing gain. Choosing $\tau$ small
keeps every actual preactivation below $2\rho$; a fixed readout
$2/\rho$ gives margin greater than one. Counting the query layer,
combination layer, $d+1$ amplifiers, and readout gives $L_d=d+4$.
The amplifier weights are still $O(d)$. This proves the teacher
resource bounds. Rounding its coefficients to $O(d\log d)$ bits
preserves the margin: the whole propagation gain is at most
$\exp(O(d\log d))$.

\subsection{Cold and fixed-gap posterior bounds}

Let $a>0$ be the hard-product exponent in the proof of
Lemma~\ref{lem:hard}. Since $L_d=d+4$, the original fixed-gain
memorizer applies. Set
\[
 q=\exp(a/100),\qquad \alpha=c_0(q-1),
\]
with the same sufficiently small constant $c_0$ as in
Lemma~\ref{lem:memorizer}. Both are fixed, and
$q^{L_d}=\exp(ad/100+O(1))$.
The relative-error sum is still $O(\alpha/(q-1))$, hence
uniformly small. The first scalar energy is at most one: its
decaying factor is
\[
 q^{-2(L_d-2)}=\exp(-ad/50+O(1)),
\]
whereas \eqref{eq:B0} gives $\log B_0\le3cd+O(\log d)$.
Choose $c<a/1000$. The witness energy satisfies
\begin{equation}\label{eq:all-derivative-memorizer}
 K_d\le q^2L_d+O(1).
\end{equation}
The cold posterior argument remains valid at each fixed
architecture. Arithmetic--geometric mean bounds its unit-margin
representative product by
\[
 \exp\{ad/100+O(1)\}\le e^{ad}
\]
for all large $d$. This is the same hard-table threshold as before.

For the actual fixed-gap scales, Lemma~\ref{lem:perturb} and
Proposition~\ref{prop:finite} apply with the energy in
\eqref{eq:all-derivative-memorizer}. Since $L_d=d+4$ and
$q^{L_d}=\exp(ad/100+O(1))$, the inverse perturbation tolerance
remains $\exp(O(d))$ without further changes. The Gaussian bulk
constant is polynomial in $d,L_d$, and the memorizer width
in~\eqref{eq:MB} is $\exp(O(d))$ uniformly over $n\le e^{cd}$.
Thus every $N\ge\exp(C_\zeta d)$ meets the finite-scale
condition~\eqref{eq:width}, for each fixed gap $\zeta=\gamma-1>0$.
The posterior product bound remains below $e^{ad}$; the hard-table
and Gibbs-to-confidence calculation gives the same failure rate
$139/1024$.

\subsection{The critical endpoint}

The modified activation satisfies the fixed-activation hypotheses
of Appendix~\ref{ac-appendix}: the same shape and weighted estimates,
bounded derivatives through order five, and the strict scalar
potential obtained by choosing $\tau$ sufficiently small. Therefore
its finite-width target-mass and canonical-circuit results apply
with constants depending only on this fixed activation.

For $L_d=d+4$ and $n\le e^{cd}$, reducing $c$ makes the depth
condition $L_d\ge C\log(d+n+L_d)$ hold for all large $d$.
A further reduction makes the posterior circuit ceiling
$\poly(d+n+L_d)$ at most
$2^{d/4}$. The width threshold
$2^{Cd}(d+n+L_d)^C$ is still at most $\exp(C'd)$.
The hard-table and confidence calculation is unchanged. This
completes the proof of Theorem~\ref{thm:all-derivatives}.

\subsection{Quantitative derivative bounds and numerical evaluation}

The distinction between bounded derivatives of each order and a
quantitative bound on their growth is important here.

\begin{proposition}\label{prop:gevrey-evaluation}
Suppose a real activation satisfies, for fixed constants $C,A\ge1$
and finite $s\ge1$,
\[
 \sup_{u\in\R}|\phi^{(j)}(u)|\le CA^j(j!)^s
 \qquad(j\ge0).
\]
On dyadic inputs of bit length $b$ in $[-R,R]$, it has nonuniform
Boolean circuits approximating its value to error $2^{-p}$, of
size polynomial in $R+1$, $p+1$, and $b+1$. The polynomial may
depend on $C,A,s$. No efficient uniform procedure for computing
the circuit's coefficients is assumed.
\end{proposition}
\begin{proof}
Replace $s$ by the integer $s_0=\lceil s\rceil$, and choose an
integer upper bound for $A$. For an integer $n$, partition the
interval into pieces of radius at most
$h=(2An^{s_0-1})^{-1}$. Taylor expansion of degree $n-1$ at the
center of a piece has remainder at most
\[
 \frac{CA^n(n!)^{s_0}}{n!}h^n
 \le C(Ahn^{s_0-1})^n=C2^{-n}.
\]
Take $n=O(p+\log C+1)$, with enough fixed slack. There are
$O((R+1)n^{s_0-1})$ pieces. Store their centers and Taylor
coefficients as nonuniform circuit constants. The coefficients
have logarithmic magnitude $O(n\log(n+1))$, and the local
argument has magnitude at most one. Rounding each coefficient
to absolute error $2^{-p}/(8n)$ controls the total coefficient
error. Polynomially many additional bits suffice for the dyadic
arithmetic, interval selection, and evaluation of each polynomial.
All sizes are polynomial in the displayed resources.
\end{proof}

This includes a global finite-order Gevrey bound. A bounded
holomorphic extension to a fixed-width strip is another sufficient
condition, by Cauchy estimates. Such conditions cannot hold for
the constructed activation: simulating its polynomial-resource
teacher at precision $O(d\log d)$ would give polynomial-size
Boolean circuits for the circuit-hard $f_d$. Thus the new variant
still contains computationally expensive information in its
activation primitive. It is not a counterexample for a standard
efficiently evaluable activation or for exact tanh.
